\documentclass[11pt]{article}

% --- packages (all standard; compiles on stock TeX Live / Overleaf) ---------
\usepackage[margin=1in]{geometry}
\usepackage[T1]{fontenc}
\usepackage{lmodern}      % scalable fonts (needed for microtype expansion on BasicTeX)
\usepackage{amsmath,amssymb,amsthm}
\usepackage{booktabs}
\usepackage{graphicx}
\usepackage{xcolor}
\usepackage{caption}
\usepackage[hidelinks]{hyperref}
\usepackage{microtype}

\graphicspath{{figures/}}

% --- theorem environments ---------------------------------------------------
\theoremstyle{plain}
\newtheorem{proposition}{Proposition}
\newtheorem{lemma}{Lemma}
\newtheorem{conjecture}{Conjecture}
\theoremstyle{definition}
\newtheorem{definition}{Definition}
\theoremstyle{remark}
\newtheorem{remark}{Remark}

% --- shorthand --------------------------------------------------------------
\newcommand{\bits}{\,\text{bits/symbol}}
\newcommand{\pictx}{\pi_{\text{ctx}}}
\newcommand{\Pglob}{P_{\text{g}}}
\newcommand{\xstar}{x^{\!*}}
\newcommand{\norm}[1]{\left\lVert #1 \right\rVert}
\newcommand{\TV}{\mathrm{TV}}
\DeclareMathOperator{\supp}{supp}

\title{\bfseries Trust Is the Attack Surface:\\[2pt]
A Cryptanalytic Reading of In-Context Memory Poisoning\\
in a Markov Mixture Model}
\author{Dan Shumow\thanks{Research-style write-up of the analysis in the
open-source project \texttt{trust-is-the-attack-surface}, built on the toy model
\texttt{markov-models-aux-context-caching}. All figures and numbers are reproduced by
\texttt{demo\_conservation.py} in that repository.}}
\date{June 2026}

\begin{document}
\maketitle

\begin{abstract}
A large language model relies on two memories: the frozen weights, and the context
window it is handed at run time. In-context learning is the model leaning on the second
memory, and it is tempting to assume a \emph{conservation law}: the more usefully a model
relies on its context, the more an attacker who poisons that context can steer it---so
that usefulness and exploitability are two names for one quantity. We test this on a
deliberately minimal model where the question has an exact answer: a global (parametric)
Markov model linearly combined with an online Markov \emph{cache} (the context window),
evaluated prequentially in bits per symbol. Reading context-poisoning as a
\emph{differential attack} makes the objects of differential cryptanalysis---per-round
difference decay, the differential hull, weak-key leverage---\emph{literally exact} here
rather than heuristic, because the predictor genuinely is a Markov chain. We find the
clean conservation law is \emph{false as stated}. Exploitability does not track
usefulness; it tracks \emph{trust}, the weight the model places on context, and trust and
usefulness coincide only under \emph{context sparsity}. With a dense (low-order) cache,
trust saturates near $0.9$ regardless of benefit, and a propagated poison stays maximally
effective \emph{exactly where the cache hurts the model by more than a bit per symbol};
the usefulness--exploitability correlation collapses from $0.97$ to $0.52$. We also exhibit
a \emph{condensation} transition: in generation a single self-reinforcing injected count is
shrugged off below a trust threshold and locks in above it. We close by stating a
falsifiable \emph{trust-saturation conjecture} for real transformers, with induction as the
prime suspect, and a defensive corollary: bound \emph{trust}, not usefulness.
\end{abstract}

% ===========================================================================
\section{Introduction}
% ===========================================================================
A language model has two memories. One is the weights---slow, expensive, frozen at
deployment, holding what the model \emph{knows}. The other is the context window---fast,
free, rewritten every turn, holding what the model has just been \emph{told}. In-context
learning is the model leaning on that second memory, and it is most of what makes these
systems feel adaptive.

Anything a model leans on, an attacker can lean on too. The natural worry is clean:
\emph{the more usefully a model relies on its context, the more an attacker who poisons
that context can steer it}. Usefulness and exploitability would then be two readings of one
quantity---the model's dependence on what it was just told. Call it the
\textbf{conservation hypothesis}. It is the kind of thing one would put in a threat model
without checking. We checked it, on a model where it has an exact answer, and it is wrong
in an instructive way.

\paragraph{Why a toy, and why a cipher.}
One cannot ask this question cleanly of a real transformer: the quantities one wants---how
far a single poisoned token's influence propagates, how much a perturbation is
amplified---are uncomputable when attention mixes content-dependently, since the mixing
operator is itself a function of the input being perturbed. So we drop to a model that
keeps the one feature under study and discards the rest: \emph{two memories, linearly
combined}. A frozen global Markov model stands in for the weights; an online Markov cache,
estimated live from the current sequence, stands in for the context window. The cache's
warm-up \emph{is} in-context learning, made of nothing but counts. This is the
from-scratch, no-neural-network version of the global-versus-in-context decomposition
studied by Bietti et al.~\cite{bietti2023}.

Reading poisoning through a cryptographic lens is natural because injecting a crafted input
to steer an output is, structurally, a \textbf{differential attack}: introduce a difference
at the input and ask how it propagates to the output. The payoff of going to a Markov model
is that the objects of differential cryptanalysis that are only ever \emph{approximations}
for a cipher become \emph{exact} here. The ``Markov cipher'' assumption that lets
cryptanalysts multiply per-round probabilities~\cite{lai1991} is normally a heuristic; here
it holds by construction, because the predictor genuinely is a Markov chain. The
worst-case per-step decay of a difference is the chain's Dobrushin coefficient, in closed
form; the differential hull is a resolvent one can actually invert; and the leverage of a
single poisoned count in a context seen $n$ times is exactly $1/(n+a)$, plain estimator
variance.

\paragraph{Contributions.}
\begin{enumerate}\itemsep2pt
\item A dictionary (\S\ref{sec:dict}) translating context-poisoning into differential
cryptanalysis, with four objects (\S\ref{sec:exact}) that are exact rather than approximate
on this architecture.
\item A measured refutation of the conservation hypothesis (\S\ref{sec:results}):
exploitability tracks \emph{trust}, not usefulness; the two coincide only under context
sparsity. With a dense cache the model is fully poisonable precisely where its context is
useless.
\item A condensation phase transition in generation (\S\ref{sec:condensation}): a single
self-reinforcing poison locks in once trust crosses a threshold.
\item A falsifiable trust-saturation conjecture for real transformers and a defensive
corollary (\S\ref{sec:conjecture}).
\end{enumerate}

% ===========================================================================
\section{The model}\label{sec:model}
% ===========================================================================
Fix a vocabulary $[V]=\{0,\dots,V-1\}$. A document is a sequence $x_{1:T}$, and a predictor
assigns $P(x_t\mid x_{1:t-1})$. We evaluate \emph{prequentially}: score the next symbol,
then update on it, and report cross-entropy in bits per symbol,
\begin{equation}\label{eq:xent}
H \;=\; -\frac1T\sum_{t=1}^{T}\log_2 P\!\left(x_t \mid x_{1:t-1}\right).
\end{equation}

\paragraph{Global model (the weights).}
The global predictor $\Pglob(\cdot\mid c)$ is a corpus-trained, frozen Markov back-off model
(interpolated Witten--Bell~\cite{goodman2001}) over orders $0..K_g$. It is a strong,
tuning-free baseline that backs unseen high-order contexts off to lower orders rather than
to the uniform distribution, so it ``knows something everywhere.'' Its internals do not
matter below; only that it is fixed at deployment.

\paragraph{Cache (the context window).}
The cache is an order-$k$ count table built online. Writing $c=(x_{t-k},\dots,x_{t-1})$ for
the current context state, let $N_c(w)$ be the number of times symbol $w$ has followed $c$
\emph{so far in this sequence}, with $n_c=\sum_w N_c(w)$. The cache starts empty and is
updated after each observation.

\paragraph{Combination (the trust knob).}
We combine the global prior with the cache evidence by a Dirichlet (Bayesian) rule: the
global model is the prior, the cache is data,
\begin{equation}\label{eq:dir}
P(w\mid c)\;=\;\frac{N_c(w)+a\,\Pglob(w\mid c)}{n_c+a}.
\end{equation}
Equivalently $P(\cdot\mid c)=\rho_c\,\widehat{P}_{\text{cache}}(\cdot\mid c)+(1-\rho_c)\,
\Pglob(\cdot\mid c)$, an adaptive linear interpolation whose weight on the cache,
\begin{equation}\label{eq:reliance}
\rho_c \;=\; \frac{n_c}{n_c+a},
\end{equation}
grows automatically with evidence. We call $\rho_c$ the \textbf{reliance} (or
\textbf{trust}): the realized weight the prediction places on context. The prior strength
$a>0$ is the only knob; small $a$ lets the cache dominate. When the cache has never seen
$c$, \eqref{eq:dir} collapses to $\Pglob$.

\paragraph{Synthetic source (a dial on how much context exists).}
Each document is generated from two layers. A single transition matrix $G$ (rows drawn from
$\mathrm{Dir}(\gamma\mathbf 1)$ with $\gamma=1$, high entropy) is shared by every document. Each
document also draws its own peaky matrix $D$ (from $\mathrm{Dir}(\delta\mathbf 1)$,
$\delta=0.05$). At each step we sample the next symbol from $D[x_{t-1}]$ with probability
$\pictx$, else from $G[x_{t-1}]$:
\begin{equation}\label{eq:source}
x_t \;\sim\; \begin{cases}D[x_{t-1}] & \text{w.p. } \pictx,\\[2pt] G[x_{t-1}] &
\text{w.p. } 1-\pictx.\end{cases}
\end{equation}
The global model, trained across many documents, can only ever learn $G$; each document's
$D$ is random and washes out in the average. So the cache's test-time gain is \emph{exactly
the in-context information the global model provably cannot have}, and $\pictx$ is a clean
dial on how much of it exists---equivalently, how much there is for an attacker to exploit.
A high-entropy $G$ is deliberate: at $\pictx=0$ documents are unstructured, so a context
recurring is \emph{not} automatic, and trust is free to vary rather than being handed out by
a small context space.

% ===========================================================================
\section{A cryptanalytic dictionary}\label{sec:dict}
% ===========================================================================
The same forward computation reads three ways. The point of the toy is that the rightmost
column is \emph{exact and computable}, where the middle column is only ever approximate.

\begin{center}\small
\begin{tabular}{p{0.31\textwidth}p{0.28\textwidth}p{0.32\textwidth}}
\toprule
Differential cryptanalysis & Transformer & This model (global $+$ cache) \\
\midrule
Plaintext difference (finite DDT) & token edit $\to$ embedding diff. & a single injected
symbol / count \\
Round function & attention $+$ MLP block & one mixture step
$M=\rho\,\text{cache}+(1{-}\rho)\,\text{global}$ \\
Markov-cipher assumption & \emph{approximation} & \textbf{literally true}---it is a Markov
chain \\
Max diff.\ probability / round & unbounded, no branch number & \textbf{Dobrushin coefficient}
$\tau(M)$, closed form \\
Differential hull (uncomputable) & sum over a trail cluster & \textbf{resolvent}
$(I-M)^{-1}$ on mean-zero space \\
Weak keys / encystment site & dead ReLUs, null tokens & \textbf{low-count contexts},
leverage $1/(n+a)$ \\
Induction head (copy + vector) & QK match-and-copy & PPM longest-match (same dual role) \\
Self-propagation & autoregressive feedback & \textbf{nonlinear P\'olya urn}, condensation \\
\bottomrule
\end{tabular}
\end{center}

Every row that is approximate for a transformer becomes closed form here because the
architecture is an affine map on the probability simplex with no softmax to hide the
difference behavior inside.

% ===========================================================================
\section{Four objects made exact}\label{sec:exact}
% ===========================================================================
A \emph{difference} is a signed measure $\delta\in\mathbb R^V$ with $\sum_w\delta(w)=0$; it
is the cryptanalyst's input difference, here a perturbation of a predictive row. Let $M$ be
the row-stochastic operator the chain applies (context state to next-symbol distribution,
lifted to a transition on context states). A difference propagates by $\delta\mapsto
\delta^{\!\top}M$.

\paragraph{(1) Per-step decay is the Dobrushin coefficient.}
\begin{definition}[Dobrushin ergodic coefficient~\cite{dobrushin1956}]
For a row-stochastic $M$,
\[
\tau(M)\;=\;\tfrac12\max_{s,s'}\norm{M(s,\cdot)-M(s',\cdot)}_1
\;=\;\max_{s,s'}\norm{M(s,\cdot)-M(s',\cdot)}_{\TV}.
\]
\end{definition}
\begin{lemma}\label{lem:dob}
For any mean-zero $\delta$, $\;\norm{\delta^{\!\top}M}_1\le \tau(M)\,\norm{\delta}_1$.
\end{lemma}
This is the literal ``maximum differential probability per round.'' AES forces it small by
design (the branch number); the mixture chain simply \emph{has} a value, computable from its
rows. The wide-trail bound that does not exist for transformers exists here.

\paragraph{(2) The differential hull is the resolvent.}
The strength of a differential attack is never one trail; it is the \emph{hull}, the sum
over every path a difference can take between fixed endpoints.
\begin{proposition}\label{prop:hull}
On a stationary stretch with fixed operator $M$ and $\tau(M)<1$, the accumulated influence
of an initial difference $\delta$ is
\[
\sum_{t\ge 0}\delta^{\!\top}M^{t}\;=\;\delta^{\!\top}(I-M)^{-1}
\qquad\text{on the mean-zero subspace,}
\]
with $\norm{\delta^{\!\top}(I-M)^{-1}}_1\le \norm{\delta}_1/(1-\tau(M))$.
\end{proposition}
The ``thick hull'' that is uncomputable for a transformer is here a Green's function---the
chain's fundamental matrix---that one can invert.
\begin{remark}
In our setting $M$ is not fixed: the cache updates as the sequence streams, so the genuine
object is a history-dependent product $\prod_t M_t$. Proposition~\ref{prop:hull} is the
stationary idealization; the \emph{measured} propagated leverage of \S\ref{sec:results} is
the honest history-dependent version, and it is still explicit because the cache update is
linear in counts.
\end{remark}

\paragraph{(3) Encystment: weak keys are low-count contexts, leverage is $1/(n+a)$.}
\begin{proposition}[Single-count leverage]\label{prop:leverage}
Fix a context $c$ with cache total $n=n_c$, and let $P=P(\xstar\mid c)$ be the current
combined probability \eqref{eq:dir} on a target $\xstar$. Injecting one additional count of
$\xstar$ into $c$ changes that probability by exactly
\begin{equation}\label{eq:lev}
\Delta P \;=\; \frac{1-P}{\,n+a+1\,}\;\le\;\frac{1}{n+a+1},
\end{equation}
independent of $\xstar$, and $\Delta P\to 0$ as $n\to\infty$.
\end{proposition}
\begin{proof}
Write $P=(N+ag)/(n+a)$ with $N=N_c(\xstar)$ and $g=\Pglob(\xstar\mid c)$. After the
injection the numerator and total each rise by one count of $\xstar$:
$P'=(N+1+ag)/(n+1+a)$. Then
\[
\Delta P=P'-P=\frac{(N+1+ag)(n+a)-(N+ag)(n+1+a)}{(n+a)(n+1+a)}
=\frac{(n+a)-(N+ag)}{(n+a)(n+1+a)}=\frac{1-P}{n+1+a},
\]
using $N+ag=(n+a)P$ in the last step.\qedhere
\end{proof}
This is textbook estimator variance: the weak key is the rare $n$-gram, and ``rare'' is a
number in the count table. Crucially, \eqref{eq:lev} \emph{shrinks} as the cache fills---a
fact that, taken alone, would mislead (\S\ref{sec:results}).

\paragraph{(4) Sporulation is a nonlinear P\'olya urn.}
In \emph{generation}, sampling a symbol increments its own count, so under high reliance the
process is positively reinforced. Reinforced (P\'olya) urns with super-linear feedback
exhibit \emph{condensation}: below a threshold the system stays diffuse; above it a single
early advantage locks in and takes over~\cite{pemantle2007,chung2003}. The order parameter is
the cache trust $\rho$, and the phenomenon is a measurable phase transition
(\S\ref{sec:condensation}).

% ===========================================================================
\section{Experimental setup}\label{sec:setup}
% ===========================================================================
We fix $V=32$, global order $K_g=2$, prior strength $a=1$, and sweep $\pictx\in\{0,0.1,
\dots,0.9\}$ for two cache orders: $k=1$ (\emph{dense}: only $32$ possible contexts, so the
cache fills regardless) and $k=3$ (\emph{sparse}: most contexts are rare). For each setting
we measure four quantities.

\begin{description}\itemsep3pt
\item[Usefulness (benefit).] Bits per symbol saved by the cache versus the global model
alone, $H_{\text{global}}-H_{\text{cache}}$ from \eqref{eq:xent}. Positive means context
helps.
\item[Trust (reliance).] The warmed-up cache weight $\rho$ of \eqref{eq:reliance}, averaged
over the held-out stream.
\item[Static leverage (control).] The one-shot, non-propagating poison of
Proposition~\ref{prop:leverage}: at each live context inject one count of the
globally-rarest symbol $\xstar=\arg\min_w\Pglob(w)$, and record the immediate swing
$P'(\xstar)-P(\xstar)$, then undo it. A thin context is always pollutable; the point is that
a one-shot poison that does not \emph{propagate} is cheap everywhere.
\item[Propagated leverage.] The honest exploitability. Inject one count of $\xstar$ at a
single context $c_0$ a quarter of the way through the document ($t_0=T/4$), then let the
cache keep streaming the \emph{true} document, and accumulate the extra probability the model
places on $\xstar$ at every later step,
\begin{equation}\label{eq:prop}
L_{\text{prop}}\;=\;\sum_{t=t_0}^{T-1}\Big(P^{\text{pois}}(\xstar\mid x_{1:t})
-P^{\text{clean}}(\xstar\mid x_{1:t})\Big).
\end{equation}
This is the discrete Green's-function influence of a unit perturbation through the
history-dependent cache operator: large exactly when the poisoned context \emph{recurs} and
is \emph{trusted}.
\end{description}

% ===========================================================================
\section{Results: exploitability tracks trust, not usefulness}\label{sec:results}
% ===========================================================================
\begin{figure}[t]
\centering
\includegraphics[width=\textwidth]{conservation_law.png}
\caption{Usefulness versus exploitability across the $\pictx$ sweep (each curve normalized to
$[0,1]$). \textbf{Left} (cache order 3, sparse / trust earned): usefulness, reliance, and
propagated leverage rise together, $\mathrm{corr}(\text{benefit},\text{propagated})=0.97$.
\textbf{Middle} (cache order 1, dense / trust saturated): reliance pins at ${\sim}0.91$ and
propagated leverage stays high \emph{even where the cache hurts by more than a bit}; the
correlation collapses to $0.52$. \textbf{Right}: the static one-shot control is
\emph{anti}-correlated with usefulness in both regimes.}
\label{fig:conservation}
\end{figure}

\paragraph{Result 1: the obvious exploitability measure is the wrong one.}
The one-shot (static) measure is what most people would reach for, and it is
\emph{anti}-correlated with usefulness ($r\approx-0.99$ at order 3, $-0.99$ at order 1). It
falls as context becomes more useful, because a warmed-up context has high counts and one
extra count barely moves anything---Proposition~\ref{prop:leverage} with large $n$. One-shot
leverage is a red herring; only the \emph{propagated} measure is the right object, exactly
the cryptographic point that a differential's value is the hull, not a single trail.

\paragraph{Result 2: dense contexts decouple trust from usefulness.}
Table~\ref{tab:order1} is the alarming regime. With only $32$ order-1 contexts the cache fills
regardless of $\pictx$, so reliance pins at ${\sim}0.91$ across the entire sweep---including
$\pictx=0$, where the cache \emph{hurts} by $1.29\bits$. Propagated leverage stays high
throughout ($1.1$--$1.3$), and $\mathrm{corr}(\text{benefit},\text{propagated})=0.52$.
\emph{The model is fully poisonable precisely where its context is useless.} Trust is handed
out for sheer quantity of data, not quality, and the linear mixture weights it by quantity.

\begin{table}[t]
\centering\small
\caption{Cache order 1 (dense / trust saturated). Reliance is flat near $0.91$ and
propagated leverage stays high even where benefit is strongly negative.}
\label{tab:order1}
\begin{tabular}{rrrrr}
\toprule
$\pictx$ & benefit (bits) & reliance & static (control) & propagated \\
\midrule
0.0 & $-1.295$ & 0.915 & 0.174 & 1.177 \\
0.1 & $-1.257$ & 0.913 & 0.174 & 1.138 \\
0.2 & $-1.097$ & 0.915 & 0.174 & 1.161 \\
0.3 & $-0.815$ & 0.914 & 0.172 & 1.145 \\
0.4 & $-0.469$ & 0.915 & 0.172 & 1.203 \\
0.5 & $-0.107$ & 0.914 & 0.170 & 1.196 \\
0.6 & $+0.230$ & 0.913 & 0.168 & 1.185 \\
0.7 & $+0.744$ & 0.916 & 0.166 & 1.110 \\
0.8 & $+1.202$ & 0.916 & 0.163 & 1.187 \\
0.9 & $+1.733$ & 0.916 & 0.159 & 1.293 \\
\bottomrule
\end{tabular}
\end{table}

\paragraph{Result 3: sparse contexts couple them.}
Table~\ref{tab:order3} is the safe regime. With many rare order-3 contexts, trust must be
\emph{earned}: a high-order context accumulates enough counts to be trusted only if it
actually recurs, and in this source recurrence is itself evidence of genuine local
structure. So reliance climbs $0.02\to0.39$ with $\pictx$, benefit climbs $-0.02\to+0.47$,
and propagated leverage climbs $0.52\to0.77$, all together:
$\mathrm{corr}(\text{benefit},\text{propagated})=0.97$. Here the conservation hypothesis is
true. This is the regime high-order / natural-language contexts plausibly inhabit---most long
prefixes are nearly unique. The repository includes a small prose calibration through the same
toy substrate, but the real-transformer version of this claim remains to be measured.

\begin{table}[t]
\centering\small
\caption{Cache order 3 (sparse / trust earned). Reliance, benefit, and propagated leverage
rise together.}
\label{tab:order3}
\begin{tabular}{rrrrr}
\toprule
$\pictx$ & benefit (bits) & reliance & static (control) & propagated \\
\midrule
0.0 & $-0.018$ & 0.020 & 0.496 & 0.522 \\
0.1 & $-0.018$ & 0.024 & 0.496 & 0.500 \\
0.2 & $-0.019$ & 0.022 & 0.495 & 0.533 \\
0.3 & $-0.017$ & 0.037 & 0.492 & 0.532 \\
0.4 & $-0.034$ & 0.050 & 0.488 & 0.504 \\
0.5 & $-0.025$ & 0.084 & 0.480 & 0.568 \\
0.6 & $-0.008$ & 0.128 & 0.470 & 0.515 \\
0.7 & $+0.064$ & 0.194 & 0.454 & 0.591 \\
0.8 & $+0.186$ & 0.251 & 0.435 & 0.656 \\
0.9 & $+0.468$ & 0.385 & 0.404 & 0.774 \\
\bottomrule
\end{tabular}
\end{table}

\begin{table}[t]
\centering\small
\caption{Per-order Pearson correlations across the $\pictx$ sweep. The law holds (propagated
tracks benefit) only in the sparse regime; the static control is anti-correlated in both.}
\label{tab:corr}
\begin{tabular}{lccc}
\toprule
cache order & corr(benefit, propagated) & corr(benefit, reliance) & corr(benefit, static) \\
\midrule
1 (dense)  & $+0.52$ & $+0.56$ & $-0.99$ \\
3 (sparse) & $+0.97$ & $+0.93$ & $-0.94$ \\
\bottomrule
\end{tabular}
\end{table}

\paragraph{Summary.}
The conservation law is \emph{real but conditional}, and the condition is context sparsity.
Where contexts are dense and trust is handed out for mere quantity of data, trust decouples
from usefulness and the model is exploitable independent of benefit. Where contexts are
sparse and trust must be earned from real recurrence, trust and usefulness move together and
exploitability is, in effect, paid for by capability. The same architecture sits in both
regimes depending only on how sparse its contexts are.

% ===========================================================================
\section{Condensation: a phase transition in generation}\label{sec:condensation}
% ===========================================================================
\begin{figure}[t]
\centering
\includegraphics[width=0.62\textwidth]{condensation.png}
\caption{Condensation. At fixed high $\pictx$, with one seeded self-loop $\xstar\to\xstar$,
poison takeover of generation (seeded frequency minus baseline) stays near zero up to
reliance ${\approx}0.9$, then turns sharply upward---the nonlinear P\'olya-urn knee.}
\label{fig:condensation}
\end{figure}

In generation the model samples a symbol and conditions on what it just produced. Under high
trust this loop is positively reinforcing, a nonlinear P\'olya urn. We seed a single
self-reinforcing poison ($\xstar\to\xstar$ a few times) at fixed $\pictx$ and sweep the prior
strength $a$, hence the reliance $\rho$, measuring how much $\xstar$ takes over the generated
output relative to an unseeded baseline (Table~\ref{tab:cond}, Fig.~\ref{fig:condensation}).
Takeover stays near zero up to $\rho\approx0.9$, then turns sharply upward: $+0.01$ at
$\rho=0.65$, $+0.15$ at $\rho=0.96$, and above $0.6$ once $\rho$ crosses $0.99$. A single
injected count is shrugged off when trust is moderate and \emph{locks in} once trust is high
enough. Poisoning in-context memory is not necessarily graded; above the knee it is a phase
change.

\begin{table}[h]
\centering\small
\caption{Condensation sweep at fixed $\pictx$: poison takeover versus warmed-up reliance on
the seeded context. Small $a$ means the cache dominates, raising reliance.}
\label{tab:cond}
\begin{tabular}{rrr}
\toprule
$a$ & reliance $\rho(\xstar)$ & poison takeover \\
\midrule
8.00 & 0.647 & $+0.011$ \\
4.00 & 0.849 & $+0.053$ \\
2.00 & 0.888 & $+0.059$ \\
1.00 & 0.958 & $+0.146$ \\
0.50 & 0.972 & $+0.373$ \\
0.25 & 0.990 & $+0.612$ \\
0.12 & 0.995 & $+0.653$ \\
0.06 & 0.999 & $+0.649$ \\
\bottomrule
\end{tabular}
\end{table}

% ===========================================================================
\section{The trust-saturation conjecture}\label{sec:conjecture}
% ===========================================================================
The toy is a toy, but it makes a sharp enough claim to name the analogous statement for real
transformers precisely, so it can be attacked.

\begin{conjecture}[Trust saturation]
A transformer's in-context exploitability---how much poisoning the context window steers
downstream generation---is governed by how much the model \emph{trusts} the relevant context
(the effective weight its computation places on it), not by how much that context helps the
task. Consequently: \textup{(a)} where context-reliance is earned from genuine, sparse
structure, as in most natural-language prediction, trust and usefulness coincide and
poisoning is roughly bounded by benefit; but \textup{(b)} there exist context regimes where
trust saturates independent of usefulness, and there a model is maximally poisonable exactly
where its context contributes least, with a condensation-like threshold above which a single
injected element can take over generation.
\end{conjecture}

There is a real reason to think transformers may often live in the safe regime (a):
high-order natural-language contexts are sparse---most long prefixes are nearly unique---which
is the earned-trust regime. The repository includes a small sample-prose calibration through
the toy substrate, but this is not yet a transformer measurement. If that is the whole story,
in-context exploitability is paid for by in-context capability.

The conjecture's bite is regime (b): that trust-saturation regimes \emph{exist and are
findable}, and are exactly the dangerous ones. The prime suspect is \textbf{induction}---the
copy mechanism~\cite{olsson2022}. An induction head trusts a context span because it
\emph{matches} a previous span, not because copying it helps; the toy's longest-match
predictor is a deliberately simple surrogate for that copy behavior, and it earns trust
on grounds orthogonal to benefit. So the
conjecture predicts that induction-dominated, highly repetitive, or heavily templated
contexts---repeated boilerplate, structured formats, long verbatim quotes, retrieval dumps
with duplicated passages---are where trust most decouples from usefulness, and where a single
planted token is most likely to undergo lock-in, with poisonability there roughly
\emph{independent} of how much the repeated content actually informs the task.

\paragraph{Experiments.}
Build a transformer analogue of propagated leverage (influence of a poisoned context token on
downstream generation, via attention-based influence or activation patching) and measure it
against an analogue of trust (effective attention weight on that token) and against task
benefit, swept from natural prose to maximally repetitive input. The conjecture predicts
leverage tracks the trust axis and not the benefit axis, with the gap opening up toward
induction-dominated inputs. Probe the condensation threshold directly: plant one adversarial
token in a repetitive context and measure takeover as a function of attention weight or
sampling temperature, looking for the knee.

\paragraph{Defensive corollary.}
If you are hardening a model against context poisoning, the quantity to watch and to bound is
\emph{trust, not usefulness}---specifically trust saturation in dense or repetitive contexts,
because a context can be fully trusted while contributing nothing, and that is the exact
configuration an attacker wants. A safe-by-default policy caps how much weight any context
span can earn on the basis of sheer recurrence, independent of whether that span is pulling
its weight on the task. Concretely on the toy: raising $a$ or otherwise capping $\rho$ bounds
the single-count leverage \eqref{eq:lev} and pulls generation back below the condensation
knee. Tuning defenses against measured usefulness would protect precisely the wrong regime.

% ===========================================================================
\section{What is claimed, and what is not}\label{sec:honesty}
% ===========================================================================
The conservation hypothesis was falsified by its own experiment, which we take to be the most
trustworthy route to the replacement. What survives is narrower: exploitability follows
trust; trust and usefulness coincide only under context sparsity; and dense-context trust
saturation is a poisonability regime hiding in plain sight. The exact pieces---Dobrushin
coefficient as per-step contraction, fundamental matrix as hull, leverage $1/(n+a)$, the
reinforced-urn condensation---are standard Markov-chain results; the contribution is the
\emph{assembly} lining them up one-to-one with differential cryptanalysis and reading the
cache as a poisonable membrane. All of it is established on the toy, exactly. None of it is
yet measured on a transformer: the conjecture of \S\ref{sec:conjecture} is stated so that it
can be attacked, and the claim that natural language sits in the earned-trust regime is
motivated by sparsity and by the repo's sample-prose calibration, but remains an analogy across
architectures until real-transformer measurements are run.

% ===========================================================================
\begin{thebibliography}{9}
\bibitem{bietti2023} A.\ Bietti, V.\ Cabannes, D.\ Bouchacourt, H.\ J\'egou, L.\ Bottou.
\emph{Birth of a Transformer: A Memory Viewpoint.} NeurIPS, 2023.
\bibitem{olsson2022} C.\ Olsson et al. \emph{In-context Learning and Induction Heads.}
Transformer Circuits Thread, Anthropic, 2022.
\bibitem{lai1991} X.\ Lai, J.\ L.\ Massey, S.\ Murphy. \emph{Markov Ciphers and Differential
Cryptanalysis.} EUROCRYPT, 1991.
\bibitem{dobrushin1956} R.\ L.\ Dobrushin. \emph{Central Limit Theorem for Nonstationary
Markov Chains.} Theory of Probability and Its Applications, 1(1):65--80, 1956.
\bibitem{goodman2001} J.\ T.\ Goodman. \emph{A Bit of Progress in Language Modeling.}
Computer Speech \& Language, 15(4):403--434, 2001.
\bibitem{pemantle2007} R.\ Pemantle. \emph{A Survey of Random Processes with Reinforcement.}
Probability Surveys, 4:1--79, 2007.
\bibitem{chung2003} F.\ Chung, S.\ Handjani, D.\ Jungreis. \emph{Generalizations of
Polya's Urn Problem.} Annals of Combinatorics, 7:141--153, 2003.
\end{thebibliography}

\end{document}
